\begin{proposition}
\label{proposition-what-exact}
\begin{reference}
\cite[Corollary 1]{WhatExact}
\end{reference}
In Situation \ref{situation-complex}, suppose $R$ is
a local Noetherian ring. The following are equivalent
\begin{enumerate}
\item $0 \to R^{n_e} \to R^{n_{e-1}} \to \ldots \to R^{n_0}$
is exact at $R^{n_e}, \ldots, R^{n_1}$, and
\item for all $i$, $1 \leq i \leq e$
the following two conditions are satisfied:
\begin{enumerate}
\item $\text{rank}(\varphi_i) = r_i$ where
$r_i = n_i - n_{i + 1} + \ldots + (-1)^{e-i-1} n_{e-1} + (-1)^{e-i} n_e$,
\item $I(\varphi_i) = R$, or $I(\varphi_i)$ contains a
regular sequence of length $i$.
\end{enumerate}
\end{enumerate}
\end{proposition}

\begin{proof}
If for some $i$ some matrix coefficient of $\varphi_i$
is not in $\mathfrak m$, then we apply Lemma \ref{lemma-add-trivial-complex}.
It is easy to see that the proposition for a complex and
for the same complex with a trivial complex added to it
are equivalent. Thus we may assume that all matrix entries
of each $\varphi_i$ are elements of the maximal ideal.
We may also assume that $e \geq 1$.

\medskip\noindent
Assume the complex is exact at $R^{n_e}, \ldots, R^{n_1}$.
Let $\mathfrak q \in \text{Ass}(R)$.
Note that the ring $R_{\mathfrak q}$ has depth $0$
and that the complex remains exact after localization at $\mathfrak q$.
We apply Lemmas \ref{lemma-exact-depth-zero-local} and
\ref{lemma-trivial-case-exact} to the localized complex
over $R_{\mathfrak q}$. We conclude that
$\varphi_{i, \mathfrak q}$ has rank $r_i$ for all $i$.
Since $R \to \bigoplus_{\mathfrak q \in \text{Ass}(R)} R_\mathfrak q$
is injective (Lemma \ref{lemma-zero-at-ass-zero}), we conclude that
$\varphi_i$ has rank $r_i$ over $R$ by the definition of rank as given
in Definition \ref{definition-rank}. Therefore we see that
$I(\varphi_i)_\mathfrak q = I(\varphi_{i, \mathfrak q})$
as the ranks do not change. Since all of the ideals
$I(\varphi_i)_{\mathfrak q}$, $e \geq i \geq 1$
are equal to $R_{\mathfrak q}$ (by the lemmas referenced above)
we conclude none of the ideals $I(\varphi_i)$ is contained in $\mathfrak q$.
This implies that $I(\varphi_e)I(\varphi_{e-1})\ldots I(\varphi_1)$
is not contained in any of the associated primes
of $R$. By Lemma \ref{lemma-silly} we may choose
$x \in I(\varphi_e)I(\varphi_{e - 1})\ldots I(\varphi_1)$,
$x \not \in \mathfrak q$ for all $\mathfrak q \in \text{Ass}(R)$.
Observe that $x$ is a nonzerodivisor (Lemma \ref{lemma-ass-zero-divisors}).
According to Lemma \ref{lemma-div-x-exact-one-less}
the complex $0 \to (R/xR)^{n_e} \to \ldots \to (R/xR)^{n_1}$ is exact
at $(R/xR)^{n_e}, \ldots, (R/xR)^{n_2}$. By induction
on $e$ all the ideals $I(\varphi_i)/xR$ have a regular
sequence of length $i - 1$. This proves that $I(\varphi_i)$
contains a regular sequence of length $i$.

\medskip\noindent
Assume (2)(a) and (2)(b) hold. We will prove that (1) holds by
induction on $\dim(R)$. If $\dim(R) = 0$, then we must have
$I(\varphi_i) = R$ for $1 \leq i \leq e$ by (2)(b). Since the
coefficients of $\varphi_i$ are contained in the maximal ideal
this can happen only if $r_i = 0$ for all $i$. By (2)(a) we
conclude that $e = 0$ and (1) holds. Assume $\dim(R) > 0$.
We claim that for any prime $\mathfrak p \subset R$ conditions
(2)(a) and (2)(b) hold for the complex
$0 \to R_\mathfrak p^{n_e} \to R_\mathfrak p^{n_{e - 1}} \to \ldots \to
R_\mathfrak p^{n_0}$ with maps $\varphi_{i, \mathfrak p}$
over $R_\mathfrak p$. Namely, since $I(\varphi_i)$ contains a
nonzero divisor, the image of $I(\varphi_i)$ in $R_\mathfrak p$
is nonzero. This implies that the rank of $\varphi_{i, \mathfrak p}$
is the same as the rank of $\varphi_i$: the rank as defined above
of a matrix $\varphi$ over a ring $R$ can only drop when passing
to an $R$-algebra $R'$ and this happens if and only if $I(\varphi)$
maps to zero in $R'$. Thus (2)(a) holds. Having said this
we know that $I(\varphi_{i, \mathfrak p}) = I(\varphi_i)_\mathfrak p$
and we see that (2)(b) is preserved under localization as well.
By induction on the dimension of $R$ we may assume the complex
is exact when localized at any nonmaximal prime $\mathfrak p$ of $R$.
Thus $\Ker(\varphi_i)/\Im(\varphi_{i + 1})$ has support contained in
$\{\mathfrak m\}$ and hence if nonzero has depth $0$.
As $I(\varphi_i) \subset \mathfrak m$ for all $i$ because
of what was said in the first paragraph of the proof, we
see that (2)(b) implies $\text{depth}(R) \geq e$.
By Lemma \ref{lemma-acyclic} we see
that the complex is exact at $R^{n_e}, \ldots, R^{n_1}$
concluding the proof.
\end{proof}